

\begin{Information}
Cette partie est plus formelle que ce que nous avons étudié jusqu'à maintenant.
Si le lecteur souhaite se concentrer sur l'utilisation pratique de l'outil, l'introduction
de ce chapitre et les deux premières sections (consacrées aux instructions de base et
le « bonus stage ») sont dispensables. Si ce que nous avons présenté jusqu'à
maintenant a semblé ardu au lecteur sur un plan formel, il est également possible
de garder l'introduction et ces deux sections pour une deuxième lecture.

Les sections consacrées aux boucles sont en revanches indispensables. Les éléments plus
formels de ces sections seront signalés et sont également facultatifs.
\end{Information}


Pour chaque notion en programmation C, nous associerons la règle d'inférence qui
correspondante, la règle utilisée de calcul de plus faible précondition qui la
régit, ainsi que des exemples d'utilisation (mais pas forcément dans cet ordre et avec
plus ou moins de liaison avec l'outil). Les premiers points seront focalisés sur
la théorie plutôt sur l'utilisation, car ces actions sont réalisées par l'outil
sans que cela nécessite d'opération particulière de la part de l'utilisateur.
Au fur et à mesure, nous nous concentrerons de plus en plus sur l'outil, en
particulier quand nous aborderons la section concernant les boucles.



\levelThreeTitle{Règle d'inférence}


Une règle d'inférence est de la forme :
$$\dfrac{P_1 \quad ... \quad P_n}{C}$$
et signifie que pour assurer que la conclusion $C$ est vraie, il faut d'abord
vérifier que les prémisses $P_1$, ..., et $P_n$ sont vraies. Lorsqu'il n'y a
pas de prémisses :
$$\dfrac{}{\quad C \quad}$$
Alors, aucune condition n'est requise pour conclure que $C$ est vraie : il s'agit un axiome.


Inversement, pour prouver qu'une certaine prémisse est vraie, il peut s'avérer
nécessaire d'utiliser une autre règle d'inférence, par exemple :
$$\dfrac{\dfrac{}{\quad P_1\quad} \quad \dfrac{P_{n_1}\quad P_{n_2}}{P_n}}{C}$$


Ce mécanisme construit progressivement l'arbre de déduction de notre raisonnement.
Dans notre raisonnement, les prémisses et conclusions manipulées seront
généralement des triplets de Hoare.


\levelThreeTitle{Triplet de Hoare}


Revenons sur la notion de triplet de Hoare :
$$\{ P \}\quad  C\quad \{ Q \}$$


Nous l'avons abordé brièvement en début de tutoriel, ce triplet exprime que si,
avant l'exécution de $C$, la propriété $P$ est vraie, et si $C$ termine, alors la
propriété $Q$ est vraie. Par exemple, si nous reprenons notre programme de
calcul de la valeur absolue (légèrement modifié) :



\CodeBlockInput{c}{abs-1.c}



Ce que nous dit Hoare, est que pour prouver notre programme, les propriétés
entre accolades dans ce programme doivent être vérifiées (nous omettons
volontairement l'une des deux postconditions pour alléger la lecture) :



\CodeBlockInput{c}{abs-2.c}



Cependant, Hoare ne nous dit pas comment automatiquement la
propriété \CodeInline{P} de ce programme. Ce que nous propose Dijkstra, est donc
un moyen de calculer, à partir d'une postcondition $Q$ et d'une commande ou d'une
liste de commandes $C$, la précondition minimale assurant $Q$ après $C$. Nous
pouvons donc, dans le programme précédent, calculer la propriété \CodeInline{P}
qui permet d'obtenir les garanties voulues.



Nous présenterons tout au long de cette partie les différents cas de la
fonction $wp$ qui, à une postcondition voulue et un programme ou une instruction,
associe la plus faible précondition qui permet de l'assurer. Nous utiliserons
cette notation pour définir le calcul correspondant à une ou plusieurs instructions :



$$wp(Instruction(s), Post) := WeakestPrecondition$$



De plus la fonction $wp$ est telle qu'elle nous garantit que le triplet de Hoare :
$$\{\ wp(C,Q)\ \}\quad C\quad \{ Q \}$$
est effectivement un triplet valide.



Nous utiliserons souvent des assertions ACSL pour présenter les notions à
venir :



\begin{CodeBlock}{c}
//@ assert ma_propriete ;
\end{CodeBlock}



Ces assertions correspondent en fait à des étapes intermédiaires possibles pour
les propriétés indiquées dans nos triplets de Hoare. Nous pouvons par exemple
reprendre le programme précédent et remplacer nos commentaires par les assertions
ACSL correspondantes (nous omettons \CodeInline{P}, car sa valeur est en fait
simplement « vrai ») :



\CodeBlockInput{c}{abs-3.c}
